\writeSectionFrame{\presentationTitle}{Das Problem}
\frameWithImageOnly{\BILDERPFAD/problem1.jpg}

\note{
	Dieses Muster dient dazu, eine Grammatik für eine eigene Sprache zu entwickeln. Sinnvoll kann das sein für 
	Probleme, die sehr häufig auftauchen und die dennoch variiert werden können. Ein Klassiker dafür ist die Suche nach
	Zeichenketten in einem Text.
}

\begin{frame}{Suche nach Zeichenketten}
	\begin{exampleblock}{Methode 1}
 		public boolean enthaeltText(String suchText, String text);
 	\end{exampleblock}
 	
 	\begin{exampleblock}{Methode 2}
 		public boolean enthaeltText(String suchText, String suchText2, String text, boolean isCaseSensitive, boolean und);
 	\end{exampleblock}
\end{frame}

\note{
	Diese auf den ersten Blick einfache Aufgabenstellung kann, wenn man sie wirklich komplett durchdenkt, 
	eine ganze Menge überladene Methoden nach sich ziehen. So könnte man Groß- und Kleinschreibung beachten und 
	mehrere Suchterme mittels und / oder verbinden wollen. 
}

\begin{frame}{Beispiel für eine Grammatik}
	\begin{center}
		\myImageWithSource{0.5}{\BILDERPFAD/problem2.png}{Matthias Geihros - Entwurfsmuster}
	\end{center}
\end{frame}

\note{
	 Dies ist ein Beispiel für eine Grammatik. Das Programm soll Testdaten erzeugen, die bestimmten Regeln folgen. 
	 Ein Ausdruck kann ein Literal sein, eine Zufallszahl, eine Wiederholung oder eine Variable. Ein Literal besteht
	 aus Zeichen, eine Variable ist irgendeine Art von Variable im Kontext. Eine Zufallszahl ist eine Ziffer von 0 bis 9
	 und eine Wiederholung gibt an, dass bestimmte Ausdrücke wiederholt werden. Eine Grammatik könnte auch Operatoren wie
	 UND / ODER usw. definieren. Diese Darstellung entspricht in etwa der erweiterten Backus-Naur-Form für kontextfreie
	 Grammatiken. Sie definiert, welche Möglichkeiten ein Interpreter interpretieren können muss. 
}

\begin{frame}{Beispiel für erzeugte Daten}
	\begin{center}
		\myImageWithSource{0.4}{\BILDERPFAD/problem3.png}{Matthias Geihros - Entwurfsmuster}
	\end{center}
\end{frame}

\note{
	   Mit dieser einfachen Grammatik können wir schon recht brauchbare Testdaten erzeugen, ohne dass tatsächlich etwas spezifisches programmiert werden muss. Jetzt stellt sich nur die Frage, wie eine solche Grammatik ausgewertet werden kann.
}